% Checkpoint35 proposed source; checkpoint34 sources remain unchanged.
\statement{A.117}{Proved}{The entire residual cyclic cohomology and coherent comparison preserve the original odd coordinates}
Use the exact residual of A.116,
\[
 R=N_1\cap O=N\cap F'\cap O,\qquad
 E_R=F/R=P\times D_o,\qquad P=F^{(2)}.
\]
All kernels are the actual closed kernels. Put
\(\rho=\operatorname{Ad}(r)\) and retain the integral matrix
\(J=\left(\begin{smallmatrix}0&-1\\1&-1\end{smallmatrix}\right)\).
Then
\[
 \begin{gathered}
 T_R=C_{E_R}(r)=(E_R')^\rho,\\
 T_R\ \xrightarrow{\ \sim\ }\ H^1(C_3,R),\qquad
 h\longmapsto[w^{-1}\rho(w)],\quad w\in F',\ wR=h,\\
 T_R=(P')^\rho\times C_{\mathrm{SL}_2(\mathbb Z_3)'}(J)
       \times\prod_{p\geq5}C_{\mathrm{SL}_2(\mathbb Z_p)}(J).
 \end{gathered}
 \tag{A.117.1}
\]
The middle identification is a pointed homeomorphism.
Writing \(T_1=C_{E_1}(r)\), the natural comparison is the exact
sequence of compact groups
\[
 1\longrightarrow M_{21}^\rho\longrightarrow T_R
                 \longrightarrow T_1\longrightarrow1.
 \tag{A.117.2}
\]
Under A.113 and (A.117.1) it is exactly the natural map of the
two cohomology spaces; its basepoint fiber corresponds to
\(M_{21}^\rho\). Its further map to A.109's two-point space
is the SAME original \(q_3\) coordinate.

For \(e=d\bmod3\), put \(R_e=y^{(1-e)/2}r\) and
\(\rho_e=\operatorname{Ad}(R_e)\). Retain the ORIGINAL full
negative frame \(J_e^\alpha\), its images in both quotients,
and \(\kappa_e(t)=(J_e^\alpha)^{-1}\alpha(t)J_e^\alpha\).
Then
\[
 \begin{split}
 T_{R,e}=C_{E_R}(R_e)
   &=(P')^{\rho_e}\times T_{3,e}\times\prod_{p\geq5}T_{p,e},\\
 H^1(C_2,T_{R,e};\kappa_e)
   &\cong H^1(C_2,(P')^{\rho_e};\kappa_e)
                   \times\prod_{p\ {\rm odd}}C_2.
 \end{split}
 \tag{A.117.3}
\]
The comparison to A.112 is the identity on every odd coordinate.
For an actual lift \(h_R\) of a positive factor frame, its coherent
defect is
\[
 w_{h_R}=(J_e^\alpha)^{-1}\alpha(h_R)^{-1}vh_R.
 \tag{A.117.4}
\]
It projects to the same defect in \(B/N_1\), and
\[
 q_3(w_{h_R})=\epsilon_A(d\bmod3)I=(-1)^{\Delta(\alpha)}I.
 \tag{A.117.5}
\]
Existence of a coherent frame is exactly the full norm condition
\[
       w_{h_R}=\kappa_e(t)t^{-1}\quad\text{for some }t\in T_{R,e}.
 \tag{A.117.6}
\]
Thus passage through the pro-two kernel preserves, but does not
determine, the original odd obstruction coordinates.

\begin{proof}
\emph{A.117.1. The entire \(C_3\) space.}
A.109 proves both \((F')^\rho=1\) and \(H^1(C_3,F')=\{*\}\).
If \(h\in(E_R')^\rho\), its lift \(w\) has
\(z_w=w^{-1}\rho(w)\in R\), and its norm telescopes to one.
A new lift \(wn\), \(n\in R\), changes the defect to
\(n^{-1}z_w\rho(n)\). Conversely every \(R\)-cocycle is an
\(F'\)-coboundary and hence arises this way.
If defects of \(w_1,w_2\) are equivalent, then
\(w_1nw_2^{-1}\) is fixed in \(F'\) for an \(n\in R\).
It is one, so the quotient elements agree.
The compact closed set
\(\{w\in F':w^{-1}\rho(w)\in R\}\) and its continuous defect
map give continuity by quotient descent; no section is chosen.
The cocycle orbit space is compact Hausdorff because the compact
action has a closed orbit relation. This proves the homeomorphism.

Since \(R\leq F'\), the abelianization of \(E_R\) is
\(\widehat{\mathbb Z}^2\). The determinant-three difference of
the cyclic matrix from the identity has zero kernel there.
Thus all pure \(r\)-centralizers lie in the derived subgroup.
The exact product in A.116 gives
\[
 E_R'=P'\times\mathrm{SL}_2(\mathbb Z_3)'
                   \times\prod_{p\geq5}\mathrm{SL}_2(\mathbb Z_p).
 \tag{A.117.7}
\]
On these derived odd coordinates the original signed cyclic
action is honest conjugation by \(J\); its sign character vanishes.
The action on \(P'\) is its natural actual action.
This proves the last line of (A.117.1).
We do not replace \((P')^\rho\) by the trivial group.

\emph{A.117.2. Exact comparison through the pro-two kernel.}
The map \(E_R'\to E_1'\) has kernel \(N_1/R=M_{21}\), by
A.116. This kernel is pro-two, but need not be central.
For a \(\rho\)-fixed element downstairs, a lift \(g\) has a
\(C_3\)-cocycle \(g^{-1}\rho(g)\) in \(M_{21}\).
The full compact coprime-action lemma of A.109 supplies \(k\)
with this defect \(k^{-1}\rho(k)\).
Then \(gk^{-1}\) is fixed and lifts the specified element.
This proves (A.117.2); its kernel is visibly \(M_{21}^\rho\).
Choosing the same lift \(w\in F'\) in both boundary constructions
proves their exact compatibility.
Choosing it again for \(K_3'\) proves compatibility with \(q_3\)
and fixes the sign convention. No completion/intersection exchange
has occurred.

\emph{A.117.3. The coherent centralizer in the stronger quotient.}
The cyclic abelian matrix for \(R_e\) is the same as that for
\(r\), since the extra factor \(y^{(1-e)/2}\) acts internally.
Consequently \(T_{R,e}\leq E_R'\).
Its odd action is conjugation by the exact \(J_e\) of A.112.
This gives the first line of (A.117.3).
The original group relations for \(J_e^\alpha\) hold before
either quotient, so the construction of \(\kappa_e\) and the
cocycle calculation in A.111 apply unchanged.
The trace argument in A.112 proves inversion on each odd torus
and its binary square quotient; it is independent of whether
the dyadic coordinate is \(P\) or \(P_*\).
This proves the second line and the identity of all odd
comparison coordinates.

The map \(T_{R,e}\to T_{1,e}\) is onto with kernel
\(M_{21}^{\rho_e}\), by the same coprime \(C_3\) argument.
A.116 also lifts the specified positive factor frames:
the two order-three elements above a given quotient element
are conjugate by the pro-two kernel.
Both occurrences of \(J_e^\alpha\) in (A.117.4) are images of
the SAME ORIGINAL word, so projecting (A.117.4) gives exactly
A.111's defect, proving (A.117.5).

For an arbitrary \(t\in T_{R,e}\) a direct ordered multiplication
gives
\[
                 w_{h_Rt}=\kappa_e(t)^{-1}w_{h_R}t.
 \tag{A.117.8}
\]
Thus \(h_Rt\) is coherent precisely when (A.117.6) holds.
In particular, over a SPECIFIED coherent frame \(h_1\) in \(B/N_1\),
choose a positive lift \(h_R\). Its defect lies in
\(M_{21}^{\rho_e}\). A coherent lift of this SAME \(h_1\) exists
exactly when
\[
 w_{h_R}=\kappa_e(k)k^{-1}
                 \quad\text{for some }k\in M_{21}^{\rho_e}.
 \tag{A.117.9}
\]
We make no coprime assertion for this last action: both the
acting group and the kernel are two-primary.
In particular there is no inferred surjectivity or injectivity
on the dyadic \(C_2\)-cohomology spaces.

\emph{A.117.4. Why the original marking remains necessary.}
In (A.117.7) the element with identity \(P'\) coordinate,
\(-I\) only at three, and identity elsewhere is an actual
central involution of \(B/R\), fixed by \(\alpha\).
The argument of A.113 applies in these exact coordinates.
Replacing a QUOTIENT negative frame by its product with this
element preserves both negative-frame equations and \(\kappa_e\),
but flips the three-adic coherent coordinate.
Its class under (A.117.1) is nonbase and maps to the corresponding
class of A.113, then to \(-I\) in A.109.
It has no fixed full lift. Such a replacement is not a change of
the unique original \(A_r\).
Equations (A.117.5)--(A.117.9) concern that original marking;
they neither evaluate its sign nor solve its norm equation.
\end{proof}

\begingroup\footnotesize\raggedright
The complete residual input is A.116, from
\path{work_A/actual_relative_perfect_congruence_residual34.md}.
The new comparison and norm proof was developed read-only during
the checkpoint34 capture and is first materialized in this
checkpoint35 source. No finite producer or target census is used.
\par\endgroup
